\documentclass[11pt,a4paper]{article}
\usepackage{fontspec,amsmath,amssymb,amsthm,microtype}
\setmainfont{lmroman10-regular.otf}[BoldFont=lmroman10-bold.otf,ItalicFont=lmroman10-italic.otf,BoldItalicFont=lmroman10-bolditalic.otf]
\setmonofont{lmmono10-regular.otf}[Scale=MatchLowercase]
\usepackage[margin=25mm]{geometry}
\usepackage[unicode,hidelinks]{hyperref}
\setlength{\emergencystretch}{3em}
\setlength{\parskip}{0.4em}
\setlength{\parindent}{0pt}
\allowdisplaybreaks[2]
\raggedbottom
\title{\texorpdfstring{Compatibility exports for independent Brownian product noise}{Compatibility exports for independent Brownian product noise}}
\author{LangevinLean mathematical exposition}
\date{}
\begin{document}\maketitle
\noindent Source module: \texttt{IPLA/\allowbreak{}Product\allowbreak{}Brownian.\allowbreak{}lean}. The accompanying HTML pairs this mathematical exposition with the exact Lean statements and records the scope of the formalization.


The exported shared module pulls an independent Brownian coordinate source to the product of that source with an arbitrary initial coupling. IPLA adds only its state-coordinate adapter in \texttt{ULAAdapter.\allowbreak{}lean}.


\section{\texorpdfstring{Brownian source, product pullback and Gaussian innovations}{Brownian source, product pullback and Gaussian innovations}}


Pulling a Brownian coordinate through the measure-preserving source projection preserves its Brownian law. Define the source Euler coordinates by

\[\xi_{k,a}=\frac{B_a((k+1)h)-B_a(kh)}{\sqrt h},\qquad h>0.\]

Gaussian increment laws give unit variance, and coordinate independence plus independent disjoint increments give independence in both \(k\) and \(a\). Product independence makes the entire source path independent of initialization. The shared proofs also identify the source noise with the Brownian endpoint noise used in the integral equation, prove square integrability and produce the product Gibbs initial state. The reference-building theorems retain their explicit stationarity assumptions at this intermediate layer; the complete canonical ULA construction proves the required Gibbs law.



The \texttt{export} commands expose the already proved declarations under the IPLA namespace. Their definitions and proofs are unchanged, so the mathematical argument is the one in the imported source modules; the dependency diagram below opens those full proofs. This file itself makes no additional analytic assumption and contains no replacement proof.


\end{document}
